\documentclass{amsart}
% For dealing with references we use the comment environment
\usepackage{verbatim}
\newenvironment{reference}{\comment}{\endcomment}
%\newenvironment{reference}{}{}
\newenvironment{slogan}{\comment}{\endcomment}
\newenvironment{history}{\comment}{\endcomment}

% For commutative diagrams we use Xy-pic
\usepackage[all]{xy}

% We use 2cell for 2-commutative diagrams.
\xyoption{2cell}
\UseAllTwocells

% We use multicol for the list of chapters between chapters
\usepackage{multicol}

% This is generally recommended for better output
\usepackage{lmodern}
\usepackage[T1]{fontenc}

% For cross-file-references

% Package for hypertext links:
\usepackage{hyperref}

% For any local file, say "hello.tex" you want to link to please

% Theorem environments.
%
\theoremstyle{plain}
\newtheorem{theorem}[subsection]{Theorem}
\newtheorem{proposition}[subsection]{Proposition}
\newtheorem{lemma}[subsection]{Lemma}

\theoremstyle{definition}
\newtheorem{definition}[subsection]{Definition}
\newtheorem{example}[subsection]{Example}
\newtheorem{exercise}[subsection]{Exercise}
\newtheorem{situation}[subsection]{Situation}

\theoremstyle{remark}
\newtheorem{remark}[subsection]{Remark}
\newtheorem{remarks}[subsection]{Remarks}

\numberwithin{equation}{subsection}

% Macros
%
\def\lim{\mathop{\mathrm{lim}}\nolimits}
\def\colim{\mathop{\mathrm{colim}}\nolimits}
\def\Spec{\mathop{\mathrm{Spec}}}
\def\Hom{\mathop{\mathrm{Hom}}\nolimits}
\def\Ext{\mathop{\mathrm{Ext}}\nolimits}
\def\SheafHom{\mathop{\mathcal{H}\!\mathit{om}}\nolimits}
\def\SheafExt{\mathop{\mathcal{E}\!\mathit{xt}}\nolimits}
\def\Sch{\mathit{Sch}}
\def\Mor{\mathop{\mathrm{Mor}}\nolimits}
\def\Ob{\mathop{\mathrm{Ob}}\nolimits}
\def\Sh{\mathop{\mathit{Sh}}\nolimits}
\def\NL{\mathop{N\!L}\nolimits}
\def\CH{\mathop{\mathrm{CH}}\nolimits}
\def\proetale{{pro\text{-}\acute{e}tale}}
\def\etale{{\acute{e}tale}}
\def\QCoh{\mathit{QCoh}}
\def\Ker{\mathop{\mathrm{Ker}}}
\def\Im{\mathop{\mathrm{Im}}}
\def\Coker{\mathop{\mathrm{Coker}}}
\def\Coim{\mathop{\mathrm{Coim}}}

% Boxtimes
%
\DeclareMathSymbol{\boxtimes}{\mathbin}{AMSa}{"02}

%
% Macros for moduli stacks/spaces
%
\def\QCohstack{\mathcal{QC}\!\mathit{oh}}
\def\Cohstack{\mathcal{C}\!\mathit{oh}}
\def\Spacesstack{\mathcal{S}\!\mathit{paces}}
\def\Quotfunctor{\mathrm{Quot}}
\def\Hilbfunctor{\mathrm{Hilb}}
\def\Curvesstack{\mathcal{C}\!\mathit{urves}}
\def\Polarizedstack{\mathcal{P}\!\mathit{olarized}}
\def\Complexesstack{\mathcal{C}\!\mathit{omplexes}}
% \Pic is the operator that assigns to X its picard group, usage \Pic(X)
% \Picardstack_{X/B} denotes the Picard stack of X over B
% \Picardfunctor_{X/B} denotes the Picard functor of X over B
\def\Pic{\mathop{\mathrm{Pic}}\nolimits}
\def\Picardstack{\mathcal{P}\!\mathit{ic}}
\def\Picardfunctor{\mathrm{Pic}}
\def\Deformationcategory{\mathcal{D}\!\mathit{ef}}

\setlength{\emergencystretch}{3em}
\begin{document}
\title[Stacks Project, Tag 08HB]{391 --- Derived objects and compatible morphisms glue across elementary distinguished squares}
\maketitle
\noindent Copyright (C) 2005--2025 Johan de Jong. Stacks Project authors. Source: \href{https://stacks.math.columbia.edu/tag/08HB}{Tag 08HB}.

\noindent Editorial difficulty note (not author text). Difficulty H2: The proof must glue derived objects from a Zariski open and an etale neighbourhood that need not be an open subspace. It builds a derived fibre triangle, then uses support on the closed complement and conservativity of etale pushforward there to prove the second restriction isomorphism and the compatible-map statement..

\noindent Editorial provenance note (not author text). History: excerpt from revision \texttt{a0444\allowbreak 6e57e\allowbreak c1fbc\allowbreak 252a8\allowbreak 71afc\allowbreak ec775\allowbreak 2fb28\allowbreak 07b14}, spaces-perfect.tex, lines 1988--2074. Reference presentation was adapted; original-author context and core support blocks were retained or added. The mathematical wording is unchanged. Licensed under GNU FDL 1.2 or later; no invariant sections or cover texts. See ../licenses/COPYING. Context blocks reproduce author definitions and supporting results; they are not additional counted problems.

\begin{definition}
\label{definition-elementary-distinguished-square}
Let $S$ be a scheme. A commutative diagram
$$
\xymatrix{
U \times_W V \ar[r] \ar[d] & V \ar[d]^f \\
U \ar[r]^j & W
}
$$
of algebraic spaces over $S$ is called an {\it elementary distinguished square}
if
\begin{enumerate}
\item $U$ is an open subspace of $W$ and $j$ is the inclusion morphism,
\item $f$ is \'etale, and
\item setting $T = W \setminus U$ (with reduced induced
subspace structure) the morphism $f^{-1}(T) \to T$ is an isomorphism.
\end{enumerate}
We will indicate this by saying: ``Let $(U \subset W, f : V \to W)$
be an elementary distinguished square.''
\end{definition}

\begin{lemma}
\label{lemma-glue}
Let $S$ be a scheme. Let $(U \subset X, V \to X)$ be an elementary
distinguished square of algebraic spaces over $S$. Suppose given
\begin{enumerate}
\item an object $A$ of $D(\mathcal{O}_U)$,
\item an object $B$ of $D(\mathcal{O}_V)$, and
\item an isomorphism $c : A|_{U \times_X V} \to B|_{U \times_X V}$.
\end{enumerate}
Then there exists an object $F$ of $D(\mathcal{O}_X)$
and isomorphisms $f : F|_U \to A$, $g : F|_V \to B$ such
that $c = g|_{U \times_X V} \circ f^{-1}|_{U \times_X V}$.
Moreover, given
\begin{enumerate}
\item an object $E$ of $D(\mathcal{O}_X)$,
\item a morphism $a : A \to E|_U$ of $D(\mathcal{O}_U)$,
\item a morphism $b : B \to E|_V$ of $D(\mathcal{O}_V)$,
\end{enumerate}
such that
$$
a|_{U \times_X V}  = b|_{U \times_X V} \circ c.
$$
Then there exists a morphism $F \to E$ in $D(\mathcal{O}_X)$
whose restriction to $U$ is $a \circ f$
and whose restriction to $V$ is $b \circ g$.
\end{lemma}

\begin{proof}
Denote $j_U$, $j_V$, $j_{U \times_X V}$ the corresponding morphisms towards
$X$. Choose a distinguished triangle
$$
F \to Rj_{U, *}A \oplus Rj_{V, *}B \to
Rj_{U \times_X V, *}(B|_{U \times_X V}) \to F[1]
$$
Here the map $Rj_{V, *}B \to Rj_{U \times_X V, *}(B|_{U \times_X V})$
is the obvious one. The map
$Rj_{U, *}A \to Rj_{U \times_X V, *}(B|_{U \times_X V})$
is the composition of
$Rj_{U, *}A \to Rj_{U \times_X V, *}(A|_{U \times_X V})$
with $Rj_{U \times_X V, *}c$. Restricting to $U$ we obtain
$$
F|_U \to A \oplus (Rj_{V, *}B)|_U \to
(Rj_{U \times_X V, *}(B|_{U \times_X V}))|_U \to F|_U[1]
$$
Denote $j : U \times_X V \to U$. Compatibility of restriction and
total direct image (Lemma \href{https://stacks.math.columbia.edu/tag/08GE}{lemma restrict direct image open (Tag 08GE)})
shows that both $(Rj_{V, *}B)|_U$ and
$(Rj_{U \times_X V, *}(B|_{U \times_X V}))|_U$
are canonically isomorphic to $Rj_*(B|_{U \times_X V})$.
Hence the second arrow of the last displayed equation has
a section, and we conclude that the morphism $F|_U \to A$ is
an isomorphism.

\medskip\noindent
To see that the morphism $F|_V \to B$ is an isomorphism we will use a trick.
Namely, choose a distinguished triangle
$$
F|_V \to B \to B' \to F[1]|_V
$$
in $D(\mathcal{O}_V)$. Since $F|_U \to A$ is an isomorphism, and since
we have the isomorphism $c : A|_{U \times_X V} \to B|_{U \times_X V}$
the restriction of $F|_V \to B$ is an isomorphism over $U \times_X V$.
Thus $B'$ is supported on $j_V^{-1}(T)$ where $T = |X| \setminus |U|$.
On the other hand, there is a morphism of distinguished triangles
$$
\xymatrix{
F \ar[r] \ar[d] &
Rj_{U, *}F|_U \oplus Rj_{V, *}F|_V \ar[r] \ar[d] &
Rj_{U \times_X V, *}F|_{U \times_X V} \ar[r] \ar[d] &
F[1] \ar[d] \\
F \ar[r] &
Rj_{U, *}A \oplus Rj_{V, *}B \ar[r] &
Rj_{U \times_X V, *}(B|_{U \times_X V}) \ar[r] &
F[1]
}
$$
The all of the vertical maps in this diagram are isomorphisms, except
for the map $Rj_{V, *}F|_V \to Rj_{V, *}B$, hence that is an isomorphism too
(Derived Categories, Lemma \href{https://stacks.math.columbia.edu/tag/014A}{lemma third isomorphism triangle (Tag 014A)}).
This implies that $Rj_{V, *}B' = 0$. Hence $B' = 0$ by
Lemma \href{https://stacks.math.columbia.edu/tag/08GG}{lemma pushforward with support in open (Tag 08GG)}.

\medskip\noindent
The existence of the morphism $F \to E$ follows
from the Mayer-Vietoris sequence for $\Hom$, see
Lemma \href{https://stacks.math.columbia.edu/tag/08H9}{lemma mayer vietoris hom (Tag 08H9)}.
\end{proof}
\end{document}
